%LTeX: language=fr-FR
\documentclass{article}
\usepackage{tasks}

\usepackage{amsmath,amssymb,amsfonts}
\usepackage{graphicx}
\usepackage{amsthm}
\usepackage[french]{babel}
\usepackage{enumitem}
\usepackage{currfile}

\usepackage{tikz}
\usetikzlibrary{calc,intersections}
\usepackage{qrcode}
\usepackage{mathrsfs}
\usepackage{enumitem}
\usepackage[french]{babel}
\usepackage{tcolorbox}
\usepackage{fancyhdr}
\usepackage{multicol}
\usepackage[headsep=3mm,footskip=5mm,includefoot,includehead]{geometry}
%showframe
\usepackage[lastpage,user]{zref}

\newtheorem*{theorem*}{Théorème}

\setlength{\parindent}{0cm}

\newenvironment{clist}{
  \begin{center}
  \begin{inparaenum}}{
    \end{inparaenum}
\end{center}}

\pagestyle{empty}
\graphicspath{ {Images/} }
\input{\currfiledir ../def.tex}
\def \Chapitre {Chapitre 4 -- Généralités sur les suites}
\def \ChapitreCourt {Ch. 4 -- Suites}


\geometry{left=1.5cm, right=1.5cm, top=1cm, bottom=1cm}
\setlength{\headheight}{14.0pt}

\theoremstyle{definition}
\newtheorem{exercice}{Exercice}

\begin{document}
\pagestyle{fancy}

\fancyhf{} % clear existing header/footer entries
% We don't need to specify the O coordinate
\fancyhead[L]{\large \textbf{\Niveau}}
\fancyhead[C]{\large \textbf{\Chapitre}}

\large

\begin{paragraph}{Pré-requis}
  \begin{itemize}
    \item Généralités sur les fonctions (chapitre 1) : image, tableau de valeurs, courbe, sens de variation
    \item Fonctions affines (chapitre 3) : points alignés, signe de $a$ et sens de variation
    \item Appliquer un pourcentage d'évolution, coefficient multiplicateur (questions flash)
    \item Python : boucle \texttt{for} (questions flash) ; tableur : lire une formule (seconde)
    \item Comparer des nombres (questions flash)
  \end{itemize}
\end{paragraph}

\begin{paragraph}{Contenus}
  \begin{itemize}
    \item Suite numérique : liste de nombres numérotés ; terme $u(n)$ de rang $n$, premier terme $u(0)$ ou $u(1)$ ; notation $u_n$, introduite après $u(n)$
    \item Génération par une formule explicite : $u_n$ en fonction de $n$
    \item Génération par une relation de récurrence : premier terme et $u_{n+1}$ en fonction de $u_n$
    \item Représentation graphique : nuage de points $(n~;~u_n)$, points non reliés
    \item Modélisation : traduire une évolution (ajout constant, pourcentage, autre) par une relation de récurrence
    \item Sens de variation : suite croissante, décroissante, ni l'une ni l'autre ; signe de $u_{n+1}-u_n$ dans des cas simples
    \item Calcul de termes au tableur (démonstration au tableau : formule explicite, formule de récurrence recopiée vers le bas ; recherche d'un seuil) et en Python (boucle \texttt{for})
  \end{itemize}
\end{paragraph}

\begin{paragraph}{Objectifs}
  \begin{itemize}
    \item Modéliser une situation simple par une suite (effectifs, populations, budgets, doses de médicament),
    \item Reconnaître si une suite est donnée par une formule explicite ou par récurrence,
    \item Calculer un terme de rang donné, dans les deux modes de génération,
    \item Traduire une phrase (« chaque année, on ajoute\dots », « chaque mois, la population augmente de $3\,\%$\dots ») en relation de récurrence,
    \item Représenter les premiers termes d'une suite par un nuage de points et lire un terme sur le graphique,
    \item Conjecturer le sens de variation à partir d'un tableau de valeurs ou d'un graphique,
    \item Étudier le signe de $u(n+1)-u(n)$ quand ce calcul est simple,
    \item Lire et compléter une formule de tableur (sur une capture d'écran) ; écrire un programme Python qui calcule un terme,
    \item Déterminer le rang à partir duquel un seuil est dépassé (tableau de valeurs, tableur).
  \end{itemize}
\end{paragraph}

\end{document}
